\chapter{Research and Existing Solutions}

\section{Research Method / Approach}
This project examines a bounded infrastructure-assurance question: after a structural Terraform change, does a previously verified IAM authorization obligation remain attached to the intended logical Amazon S3 asset? The mid-stage work implements a deterministic analysis prototype and controlled scenario tests. It does not yet provide a completed comparative evaluation. Accordingly, the design and evidence below describe a research question and proposed evaluation, not a demonstrated advantage or a claim of novelty.

\section{Security Analysis of Infrastructure as Code}
Empirical work has studied whether security practices appear in IaC configurations and how practitioners use static-analysis findings. Verdet et al. analyze Terraform security-practice adoption across 812 open-source GitHub projects covering AWS, Azure, and Google Cloud. Their analysis uses configuration scanning and reports variation by practice category, including comparatively frequent access-policy practices and less frequent encryption-at-rest practices \cite{verdet2025terraform}. This establishes a broad empirical view of practice adoption in a public-repository sample. It does not study whether a specific verified authorization requirement remains associated with the same logical asset across a refactor, and its sample does not represent all deployed infrastructure.

Hu et al. characterize 491 static-analysis alerts from ten open-source Terraform repositories and one proprietary repository \cite{hu2023alerts}. They identify ten alert categories, five security-related, and examine practitioners' views of alert actionability and the relationship between resource dependencies and alert occurrence. This work shows that alert context and dependencies matter when interpreting scanner output. Its unit of analysis is the alert and practitioner response; it does not evaluate baseline-to-successor continuity of a particular IAM/S3 obligation. The repository count and one-company practitioner survey also bound how broadly its results can be generalized.

Together, these studies motivate automated IaC security checks while leaving a narrower continuity question open for this project to examine. They should not be read as evidence that existing scanners cannot address continuity: the studies' stated samples and research questions do not establish that comparison.

\section{Terraform Refactoring and Authorization Semantics}
Terraform's \texttt{moved} block records a prior and new resource or module address. Terraform checks state for the object at the prior address and updates its address before producing a plan, avoiding destruction solely due to the address change \cite{hashicorpMoved}. This provides explicit configuration-level evidence about address refactoring and state association. It does not, by itself, establish that an IAM policy's principal, actions, and resource scope still describe the intended protected asset after configuration changes.

That distinction matters for S3 because policy scope is expressed through resource identifiers: AWS documentation distinguishes bucket and object resources and describes the ARN forms used for each \cite{awsS3Policy}. A syntactically valid or locally compliant statement can still require interpretation against the intended bucket or object. Terraform address continuity, policy validity, and obligation-to-asset continuity are related but separate questions; the sources above do not claim to solve all three together.

\section{Research Gap and Proposed Evaluation}
The bounded gap investigated here is whether an analyzer can (1) verify a registered IAM/S3 obligation on a baseline snapshot, (2) identify a supported candidate successor after a Terraform structural change, and (3) evaluate whether the same obligation holds for that successor, returning uncertainty when correspondence or semantics are ambiguous. This is a project research question, not a proven gap across all prior work or a novelty claim. The internal research design defines the intended comparison and cautions that the claim should be narrowed if the method adds no useful evidence \cite{asentSableResearchDesign}.

The prototype's proposed evaluation compares the implemented B0--B4 configurations in \path{backend/sable/baselines.py} on controlled scenarios, including candidate-only checks, baseline verification, explicit move evidence, structural correspondence, and obligation projection. The scenarios can establish expected behavior for their stated fixtures; they do not substitute for an independently sampled benchmark. A defensible comparison must report its scenario construction, reference outcomes, errors and uncertainty, and limitations, and must avoid treating synthetic fixtures as live-cloud effectiveness. The current report records implementation and focused correctness evidence only: no completed comparative benchmark, independent effectiveness result, real-world error rate, or live deployment is claimed.

\section{Proposed Solution}
SABLE parses baseline and candidate snapshots, verifies the registered obligation on the baseline, ranks candidate successors using supported structural signals, evaluates the obligation on a unique supported successor, and emits \texttt{PRESERVED}, \texttt{REGRESSED}, or \texttt{UNKNOWN}. The shared workflow binds emitted evidence to the candidate snapshot and module input hash before the final gate uses it. Chapter 5 describes the implemented components; Chapter 8 presents their acceptance tests and evidence boundaries.

\section{Complex Computing Problem and Technical Challenges}
The problem combines declarative-language normalization, bounded resolution of references across local modules, successor selection among resources, and relational authorization evaluation. The analyzer must keep resource identity evidence separate from policy validity: a valid policy attached to an unrelated bucket must not be reported as preservation. Ambiguous or unsupported inputs therefore need an explicit uncertainty outcome. The implemented subset addresses HCL normalization, local-module references, structural correspondence, action/principal/resource checks, and evidence freshness and integrity; full Terraform evaluation, complete AWS policy semantics, and live cloud state remain outside scope.

\section{Research Findings and Design Decisions}
The literature establishes (a) observed adoption patterns for selected IaC security practices in public repositories, (b) alert categories and practitioner considerations for Terraform static analysis, (c) Terraform's address-refactoring behavior, and (d) AWS's distinction between policy resource scopes. It does not establish that SABLE improves on existing tools. Current local tests show that selected fixtures produce expected statuses; they do not demonstrate novelty, general effectiveness, or an independent error rate. The research contribution remains conditional on completing the planned comparison and reporting its results and limitations.
